\scanpage{47}\mspage{53}
\uline{Morphismes \deleted{\unsure{co}}homotopiquement propres, homotopiquement lisses}

Soit $f:X\to Y$ morphisme de catégories, considérons $y_0\in\mathrm{Ob}\,\unsure{Y_0|Y}$, et le foncteur canonique
\[
 \underset{\substack{\text{fibre de }X\\\text{en }y_0}}{X_{y_0}}\xrightarrow{\varphi}X_{/y_0}
 \qquad \varphi(x)=(x,\mathrm{id}_{f(x)})\tag*{(1)}
\]
Soit $(x_0,u_0)\in\mathrm{Ob}\,X_{/y_0}$ donc
\[
 \begin{aligned}x_0&\in\mathrm{Ob}\,X\\u_0&:f(x_0)\to y_0\end{aligned}
\]
considérons la catégorie coinduite par $\varphi$ sur $(x_0,u_0)$
\[
 {}_{(x_0,u_0)\backslash}(X_{y_0})\deleted{\unsure{/(x_0,u_0)}},
\]
pour tester si $\varphi$ est une \deleted{\unsure{éq}} \unsure{hotéquivalence|hotoéquivalence} \inserted{co}stricte\inserted{ i.e. \unsure{homol.|homol.g}$^{\mathrm t}$ cofinal}. \deleted{On} On va l'interpréter via le \deleted{morphisme} \inserted{foncteur} induit par $f$
\[
 {}_{x_0\backslash}X\xrightarrow{\psi}{}_{f(x_0)\backslash}Y
\]
allant de la cat. des spécialisations de $x_0$ \inserted{de $X$,} vers celle des spécialisations de $f(x_0)$ dans $Y$. On trouve alors un iso can. \unsure{avec|entre} la fibre de $\psi$ en $u_0$

\diagram{\begin{tikzcd}[ampersand replacement=\&,row sep=small]
 x_0\arrow[r,dashed]\&x\\
 f(x_0)\arrow[r,"u_0"]\&y_0
\end{tikzcd}}
\[
 \begin{aligned}
 &\text{Catégorie coinduite }{}_{(x_0,u_0)\backslash} X_{y_0}\text{ par }X_{y_0}\text{ au dessus}\\
 &\text{de \unsure{de} l'objet }(x_0,u_0)\text{ de }X_{/y_0}\\
 &\simeq\text{ Catégorie fibre }({}_{x_0\backslash}X)_{u_0}\text{ de }{}_{x_0\backslash}X\text{ au dessus}\\
 &\text{de } {}_{f(x_0)\backslash}Y\text{ en }u_0
 \end{aligned}
\]
Donc

\uline{Proposition} Pour que les foncteurs $\varphi:X_{y_0}\to X_{/y_0}$ \inserted{($y_0\in\mathrm{Ob}\,Y$)} soient des hotéquiv. costrictes, il f. et s. que les fibres des foncteurs ${}_{x_0\backslash}X\longrightarrow {}_{f(x_0)\backslash}Y$ induits par $f$ soient contractiles.
% archive sheet number: 46
